% =========================================================
% Proof: Necessary Condition for an Extremum
% Source: volume-iii/analysis/differentiation/notes/notes-extrema.tex
% Mode: standard
% =========================================================
 
\subsection*{Necessary Condition for an Extremum}
\label{prf:necessary-condition-extremum}
 
\begin{remark*}[Return]
\hyperref[cor:necessary-condition-extremum]{$\leftarrow$ Back to Corollary (Necessary Condition for an Extremum) in Notes}
\end{remark*}
 
\begin{proof}
Let $I \subseteq \mathbb{R}$ be an open interval, let $f : I \to \mathbb{R}$,
and let $c \in I$. Assume $f$ has a relative extremum at $c$. The claim is
that either $f'(c)$ does not exist, or $f'(c) = 0$.
 
If $f'(c)$ does not exist, the conclusion holds immediately. If $f'(c)$
exists, then by \ByThm{Interior Extremum Theorem}, $f'(c) = 0$.
In either case the disjunction holds.
\end{proof}
 
\begin{remark*}[Proof shape]
A two-case exhaustion on whether $f'(c)$ exists. The first case is
immediate by definition of the disjunction. The second case is a direct
application of the Interior Extremum Theorem. No calculation is involved —
the proof is pure logical structure built on one previously proved result.
 
\FlashProof{Case 1: $f'(c)$ does not exist — first disjunct holds trivially.
  Case 2: $f'(c)$ exists — Interior Extremum Theorem gives $f'(c) = 0$,
  so second disjunct holds.}
\end{remark*}
 
\begin{remark*}[Commentary]
\textbf{(a) Why is this a corollary rather than a theorem?}
The Interior Extremum Theorem already does all the mathematical work. This
corollary just packages the conclusion in a more useful logical form:
$f'(c)$ does not exist $\vee$ $f'(c) = 0$. This disjunctive form is what
gets used in practice when hunting for critical points — it tells you exactly
what to look for without the assumption that $f$ is differentiable.
 
\textbf{(b) Why does the proof have no steps?}
Two logical cases, each a single sentence. The first is definitional;
the second is a theorem citation. There is no calculation to perform and
no construction to make. Recognising when a proof requires no `\ProofStep`
machinery — just clean case logic — is part of calibrating proof structure
to actual content.
 
\textbf{(c) What does the contrapositive say?}
If $f'(c)$ exists and $f'(c) \neq 0$, then $f$ has no relative extremum
at $c$. This is the working form: compute $f'$, find where it is zero or
undefined, and only examine those points as candidates. Every other interior
point is automatically not an extremum.
\end{remark*}
 
\begin{remark*}[Dependencies]
\begin{itemize}
  \item \hyperref[thm:interior-extremum-theorem]{Interior Extremum Theorem}:
    The single result applied in Case 2. The entire proof reduces to one
    invocation of this theorem.
\end{itemize}
 
\FlashUses{thm:interior-extremum-theorem}
\end{remark*}